We deduce Rallis's multiplicity-one theorem for twisted Jacquet
modules, as recorded in Prasad~\cite[Theorem~1]{PrasadShalika}, from
Theorem~\ref{thm:main-linear}. This argument also gives the theta-transfer
construction used for the unitary similitude groups below.

Write
\[
 N_{2,2}=\left\{n(X)=\begin{pmatrix}I_2&X\\0&I_2\end{pmatrix}:
                                  X\in\operatorname{M}_2(F)\right\},
 \qquad S_4=\Delta\GL_2(F)\ltimes N_{2,2},
\]
where $\Delta g=\diag(g,g)$, and put
$\psi_N(n(X))=\psi(\tr X)$. Since trace is invariant under
conjugation, $\sigma\boxtimes\psi_N$ defines a representation
of $S_4$ for every smooth representation $\sigma$ of $\GL_2(F)$.

\begin{corollary}[Rallis; see Prasad]\label{cor:gl-shalika}
For irreducible admissible representations $\Pi$ of $\GL_4(F)$
and $\sigma$ of $\GL_2(F)$,
\[
 \dim\Hom_{S_4}(\Pi,\sigma\boxtimes\psi_N)\leq1.
\]
Equivalently, $\dim\Hom_{\GL_2(F)}(\Pi_{N_{2,2},\psi_N},\sigma)\leq1$.
\end{corollary}

We use the equal-rank type-II oscillator normalization on
$\mathcal S(\operatorname{M}_4(F))$ given by
\[
 ((g,h)\phi)(x)=|\det(g^{-1}h)|_F^2\phi(g^{-1}xh).
\]
The big theta lift $\Theta(\tau)$ is characterized by the maximal
$\tau$-isotypic quotient $\tau\boxtimes\Theta(\tau)$ of this
oscillator. We use two facts about this lift. First, with
$L=\GL_2(F)\times\GL_2(F)$ embedded block diagonally,
Gan's big-lift period identity is
\begin{equation}\label{eq:gl-big-transfer}
 \Hom_{S_4}(\Theta(\tau),\sigma\boxtimes\psi_N)
 \simeq\Hom_L(\tau^\vee,\sigma\boxtimes\mathbf1).
\end{equation}
See \cite[Theorem~3.1]{GanPeriods}; the big-lift formulation
also appears in \cite[Proposition~3.5]{YangV1}.
This is an identity for the big lift of an arbitrary irreducible
admissible $\tau$. Second, equal-rank type-II correspondence gives
\begin{equation}\label{eq:gl-theta-quotient}
 q_\Pi:\Theta(\Pi^\vee)\twoheadrightarrow\Pi
\end{equation}
for every irreducible admissible $\Pi$.
This follows from M\'{i}nguez~\cite[Theorem~1]{Minguez08};
for the oscillator normalization above it is recalled in
\cite[Theorem~1.2 and the preceding paragraph]{FangSunXue}.

\begin{proof}
Precomposition with the surjection \eqref{eq:gl-theta-quotient},
followed by \eqref{eq:gl-big-transfer} with $\tau=\Pi^\vee$, gives
\begin{equation}\label{eq:gl-precomposition}
 \Hom_{S_4}(\Pi,\sigma\boxtimes\psi_N)
 \lhook\joinrel\longrightarrow
 \Hom_L(\Pi,\sigma\boxtimes\mathbf1).
\end{equation}
Here admissibility identifies $(\Pi^\vee)^\vee$ with $\Pi$.
To place the coefficient on the second block, let
$w=\left(\begin{smallmatrix}0&I_2\\I_2&0\end{smallmatrix}\right)$.
The block swap and evaluation pairing give an injection
\[
 \begin{aligned}
 \Hom_L(\Pi,\sigma\boxtimes\mathbf1)
 &\lhook\joinrel\longrightarrow
 \Hom_{H_{\mathrm L}}(\Pi\boxtimes\sigma^\vee,\mathbb C),\\
 A&\longmapsto
 \bigl[v\otimes\eta\longmapsto
          \langle A(\Pi(w)v),\eta\rangle\bigr].
 \end{aligned}
\]
Indeed, $w\diag(a,b)=\diag(b,a)w$ and the evaluation pairing is
$\GL_2(F)$-invariant. Injectivity follows because the smooth dual
of $\sigma$ separates vectors and $\Pi(w)$ is invertible.
Theorem~\ref{thm:main-linear} bounds the last Hom space by one,
since $\Pi\boxtimes\sigma^\vee$ is irreducible admissible
by~\cite[Proposition~2.16]{BZ76}.
\end{proof}


